\chapter{Conclusiones y trabajos futuros}
\label{ch:conclusiones}

% Este capítulo cubre el apartado "Conclusiones" de la guía: logros obtenidos y
% sugerencias de mejoras o ampliaciones futuras.

\section{Conclusiones}
\label{sec:conclusiones}

Harmony Hub ha permitido materializar una solución funcional para un problema
que suele quedar desatendido en los servicios musicales generalistas: traducir
un momento concreto del usuario a una propuesta musical con intención
funcional, trazabilidad y capacidad de adaptación progresiva. El resultado no
es únicamente una aplicación móvil conectada con Spotify, sino una arquitectura
híbrida en la que conviven contexto, reglas explícitas, catálogo musical
estructurado, aprendizaje individual del perfil y un clasificador supervisado
global deliberadamente prudente.

Desde el punto de vista técnico, el trabajo demuestra que es viable construir
una experiencia completa a partir de un flujo breve de check-in, una capa de
recomendación explicable y una materialización final en una plataforma real de
escucha. Desde el punto de vista metodológico, el proyecto aporta una
implementación donde el aprendizaje no se fuerza desde el principio, sino que
se reparte en dos niveles: una capa individual que ya puede madurar y hacerse
visible en la experiencia del usuario, y una capa supervisada global que solo
entra cuando existe evidencia suficiente y que deja trazas auditables de
cuándo actúa, cuándo se abstiene y por qué.

La aportación principal del TFG no reside en competir con motores comerciales a
gran escala, sino en integrar en una misma solución móvil la dimensión
contextual, el uso cotidiano de la música para regularse y un aprendizaje
progresivo que puede justificarse académicamente.

El estado actual del proyecto refuerza además una idea importante para su
defensa: el resultado alcanzado no es un sistema bloqueado por falta de ML,
sino una aplicación que ya personaliza de forma observable mediante
aprendizaje individual del perfil y que activa la capa supervisada global
cuando la evidencia acumulada y la confianza local de la sesión lo justifican.
La trazabilidad del sistema permite mostrar no solo cuándo interviene el
modelo, sino también con qué probabilidad, con qué umbral y con qué peso
respecto a la heurística base.

\section{Cumplimiento de objetivos}
\label{sec:cumplimiento_objetivos}

Los objetivos planteados al inicio del trabajo pueden considerarse cumplidos en
un grado suficiente para una entrega funcional y defendible. La Tabla
\ref{tab:cumplimiento_objetivos} resume esta evaluación.

\begin{table}[H]
\centering
\scriptsize
\renewcommand{\arraystretch}{1.12}
\begin{tabularx}{\textwidth}{|P{3.9cm}|Y|P{2.6cm}|}
\hline
\textbf{Objetivo} & \textbf{Síntesis de cumplimiento} & \textbf{Estado} \\
\hline
Modelos de aprendizaje automático para relacionar estado emocional y rasgos musicales. & Se implementó un enfoque híbrido con heurísticas, aprendizaje individual del perfil y un clasificador supervisado de sesión con regresión logística. No se usaron \textit{reinforcement learning} ni clustering como solución final, pero sí un modelo trazable y defendible. & Cumplido con enfoque alternativo \\
\hline
Interfaz móvil intuitiva para evaluar estado de ánimo y visualizar feedback. & La app incluye acceso, autenticación, check-in guiado, medición opcional del entorno, recomendación, feedback e histórico básico. Las pruebas funcionales y el manual muestran un flujo usable de extremo a extremo. & Cumplido \\
\hline
Integración de servicios en la nube mediante API REST. & El backend en FastAPI expone endpoints para autenticación con Spotify, check-ins, recomendaciones, playlists y feedback. Firestore persiste sesiones, catálogo y señales de aprendizaje, y Spotify materializa el resultado final. & Cumplido \\
\hline
Procesamiento del ruido ambiental con APIs del dispositivo móvil. & La app mide opcionalmente el entorno acústico y el backend clasifica esa señal en categorías funcionales que pueden modificar el perfil objetivo y las consultas auxiliares. & Cumplido en el alcance del TFG \\
\hline
Seguridad en datos sensibles y comunicación app-backend. & El sistema integra autenticación de usuarios, separación entre datos públicos y privados, reglas de seguridad de Firestore, control de acceso y gestión de credenciales fuera del repositorio. & Cumplido \\
\hline
\end{tabularx}
\caption{Cumplimiento de los objetivos principales del TFG.}
\label{tab:cumplimiento_objetivos}
\end{table}


\section{Limitaciones}
\label{sec:limitaciones}

El proyecto presenta varias limitaciones que conviene explicitar.

En primer lugar, la integración con Spotify introduce una dependencia externa
real. Aunque el sistema ya minimiza esa fragilidad apoyándose en el catálogo
\texttt{msd\_tracks}, la materialización final de playlists sigue condicionada
por permisos, disponibilidad de endpoints y posibles respuestas de
rate limit. Esta limitación existe porque el proyecto necesita cerrar el ciclo
en una plataforma musical real y, por tanto, no controla por completo la capa
de ejecución final.

En segundo lugar, la validación longitudinal del aprendizaje se ha centrado en
un único perfil de usuario. Esta decisión ha sido útil para observar mejor la
evolución temporal del sistema, pero limita la generalización experimental de
los resultados en términos poblacionales. Se asumió este enfoque para
priorizar la trazabilidad de la adaptación y evitar mezclar demasiadas
variables humanas en una validación de alcance TFG.

En tercer lugar, el conjunto supervisado actual ya no es pequeño para un TFG
de este alcance, pero sigue estando lejos de un escenario multiusuario amplio.
Los metadatos finales del modelo recogen 66 ejemplos de sesión, con 54 útiles
frente a 12 no útiles. Esta base resulta suficiente para activar el
clasificador dentro del sistema entregado, pero no equivale todavía a una
validación exhaustiva en producción ni a una cobertura poblacional amplia. La
limitación se explica por el coste real de obtener sesiones cerradas con
feedback honesto y por la decisión metodológica de no inflar el dataset con
datos artificiales.

Por último, Harmony Hub debe entenderse como un prototipo funcional orientado a
bienestar cotidiano y regulación ligera, no como una herramienta clínica ni
como un sistema de prescripción terapéutica. Esta frontera es deliberada:
permite mantener una propuesta técnicamente defendible, respetuosa con el dato
emocional y coherente con las capacidades reales del proyecto.

\section{Trabajos futuros}
\label{sec:trabajos_futuros}

Las líneas de evolución más interesantes del proyecto se concentran en cuatro
frentes.

La primera consiste en ampliar y equilibrar el conjunto de datos supervisados.
El siguiente salto de calidad ya no depende tanto de acumular sesiones útiles,
sino de recoger más contraejemplos honestos y mayor diversidad de resultados
negativos o ambiguos. Solo así podrá evaluarse el clasificador global con una
base empírica más sólida y comprobar mejor su comportamiento fuera del perfil
único analizado.

La segunda línea pasa por enriquecer el catálogo musical estructurado. Aunque
el salto dado con \texttt{msd\_tracks} ya mejora mucho la reproducibilidad del
ranking, sería valioso ampliar la base de datos con más canciones, géneros,
perfiles funcionales y descriptores semánticos. Ese crecimiento permitiría
profesionalizar el procesamiento de datos, revisar la calidad del etiquetado y
reducir todavía más la dependencia de búsquedas externas o de materializaciones
frágiles.

La tercera se relaciona con la explicabilidad. El sistema ya dispone de una
base interpretable con regresión logística, metadatos, auditoría y model
card, pero podría avanzar hacia visualizaciones más ricas del porqué de cada
recomendación y de cómo se combinan heurística, aprendizaje individual y
clasificador supervisado.

La cuarta línea se orienta a producto y evaluación. Sería relevante validar la
aplicación con más usuarios, observar patrones de uso menos homogéneos y
estudiar si ciertas configuraciones favorecen mayor utilidad percibida, mayor
adherencia o una experiencia más clara para perfiles de uso distintos. En ese
escenario tendría sentido abrir una colaboración piloto con agentes reales del
ámbito del bienestar y la salud mental digital, por ejemplo plataformas como
\textit{ifeel}, \textit{TherapyChat} o \textit{Headspace Health}, con el fin de
contrastar la utilidad del sistema en contextos de acompañamiento no clínico,
mejorar la calidad de la base de datos y someter el producto a criterios más
exigentes de profesionalización.

En conjunto, el trabajo deja una base sólida sobre la que seguir creciendo. La
aplicación ya demuestra que la música puede tratarse como una herramienta de
acompañamiento contextual con trazabilidad, prudencia técnica y capacidad real
de aprendizaje. El recorrido futuro no consiste en rehacer lo construido, sino
en ampliarlo con más datos, más contraste y más colaboración para acercarlo a
un uso cada vez más útil, responsable y humano.
